\section{The good auxiliary label on an excitation component}
\label{sec:replica_good_auxiliary_label_bound}

A physical excitation measurement acts on a different tensor factor from
an auxiliary whole-label projection. The selected component consequently
retains that label, and the compressed good-copy bound applies to it.

\begin{theorem}[Good auxiliary bound on a physical component]
\label{thm:replica_good_auxiliary_label_bound}
\lean{Matrix.replicaExcitationComponent_goodAuxiliary_labelEntropy_lower}
\leanok
Let $A,C,D$ be finite coordinate sets, let $\Omega\in\mathbb C^A$, and
let $B\subseteq\{1,\ldots,k\}$ have size $r$. Write $R_B$ for the literal
excitation operator, with factor $I-|\Omega\rangle\langle\Omega|$ on $B$
and $|\Omega\rangle\langle\Omega|$ on its complement. Enumerate the
complement first and $B$ second to define the good-copy auxiliary observable
$F_g$. Let $P_\lambda$ be the actual whole-copy central projection on
$\mathbb C^{C^k}$, with irreducible dimension $d_\lambda$, and put
\[
 \tau=\log d_\lambda-r\log|C|-\log\binom{k}{r}.
\]
For $v\in\mathbb C^{A^k}\otimes\mathbb C^{C^k}\otimes\mathbb C^D$
satisfying $(I\otimes P_\lambda\otimes I_D)v=v$, define
$w=(R_B\otimes I)v$. Then
\[
 \tau\operatorname{Re}\langle w,w\rangle
 \le\operatorname{Re}\langle w,(I\otimes F_g\otimes I_D)w\rangle.
\]
No normalization is required. Empty coordinate sets, zero copies and zero
components are included, with the total logarithm convention $\log0=0$.
\end{theorem}
\begin{proof}
\leanok
\uses{thm:good_auxiliary_label_compression,thm:replica_excitation_preserves_fixed}
The identity physical permutation and the auxiliary operator
$P_\lambda\otimes I_D$ give
$(I\otimes P_\lambda\otimes I_D)w=w$. The coordinate split has exactly
$k-r$ good and $r$ bad copies. Lift the positive semidefinite difference
$P_\lambda F_gP_\lambda-\tau P_\lambda$ by the identity on the physical
and additional auxiliary registers. Its quadratic form on $w$ is
nonnegative. The Hermitian label projection fixes $w$, and hence cancels
on both sides of that quadratic form, giving the required inequality.
\end{proof}

For a unit ground vector, $R_B$ is the actual orthogonal excitation
projection. This proves the label-support contribution to
\cite[Section~7, equation \textup{(comparator:good-auxiliary)}]{OpenAI2026AreaLaw}.
The numerical whole-label threshold and metric comparison are separate
steps, and no commutation with a band metric is used.
